\chapter{Future Work}\label{ch:futurework}

The experimental work of Chapters~\ref{ch:methods}--\ref{ch:results} characterizes the plant: the actuators and the joints they drive. The remaining task of the dissertation project is the controller --- a synthetic nervous system that schedules those actuators the way the nervous system schedules muscles --- and eventually the physical robot that the controller runs on. This chapter lays out the planned work in four stages: a neuromechanical simulation pipeline (Section~\ref{sec:fw_pipeline}), a two-layer CPG with full proprioceptive and contact feedback (Section~\ref{sec:fw_cpg}), extension of the sensorium with vestibular, ocular, and cerebellar contributions (Section~\ref{sec:fw_senses}), and physical realization on a treadmill robot with real-time neural control (Section~\ref{sec:fw_hardware}). Some of the work described here is planned for the remainder of the dissertation; the later stages extend naturally into postdoctoral research.

\section{A Neuromechanical Simulation Pipeline}\label{sec:fw_pipeline}

The first stage converts the validated biomechanical design into a physics simulation that a neural controller can drive. The pipeline has three parts.

\subsection{From OpenSim to MuJoCo}

The starting point is the author's modified OpenSim model: the Gait2392-derived robot variant in which the actuator set has been reduced and adapted to the BPA hardware and the muscle origin, insertion, and wrapping points have been updated to the optimized robot routing \citep{bolen_determination_2019, morrow_optimization_2020}. The model will be kept in OpenSim 4.x format and converted to MuJoCo \citep{todorov_mujoco_2012} using MyoConverter \citep{caggiano_myosuite_2022, wang_myosim_2022}, which converts OpenSim musculoskeletal models to MuJoCo models with optimized muscle kinematics and kinetics. The conversion will be sanity-checked by inspecting the generated model, verifying joint naming and ordering for later sensor and actuator mapping, confirming that the expected number of muscle-tendon actuators exists with plausible attachment geometry, and fixing any conversion gaps by revising the OpenSim model or accepting manual edits in the MuJoCo model.

The MuJoCo model will represent actuation with muscle-like actuators that carry an internal activation state, so that SNS outputs can drive them through activation dynamics. Each timestep, MuJoCo updates the kinematics and the muscle variables (e.g., muscle-tendon lengths and velocities) needed to compute proprioceptive signals.

\subsection{Coupling the Synthetic Nervous System}

The neural controller will be implemented with SNS-Toolbox \citep{nourse_sns_2023}, which supports large networks of spiking and non-spiking neurons and couples directly to MuJoCo. At each timestep, SNS receives Ia, Ib, and II proprioceptive channels plus binary heel and toe contact signals from the MuJoCo model, processes them through the neural core described below, and returns per-muscle activation commands that are fed into the MuJoCo muscle activation states.

The coupling will be brought up incrementally. Before any gains are tuned against human gait, the mechanics-only loop will be verified with open-loop activation patterns, confirming the sign and magnitude behavior of torques and forces. The SNS will then be added with a minimal sensor set --- contact signals plus a single proprioceptive channel --- and only later expanded to the full Ia/Ib/II set. Once stable, the activation strengths and neuron-to-muscle mapping parameters will be tuned across a few walking speeds. A Simscape-based mechanical plant could substitute for MuJoCo for validation of specific mechanical subsystems, but Simscape does not natively host the spiking and non-spiking neuron models of an SNS, so MuJoCo remains the primary platform.

\section{A Two-Layer CPG with Sensory Feedback}\label{sec:fw_cpg}

\subsection{Architecture}

The neural core follows the two-level organization of the mammalian locomotor CPG established by the Rybak group \citep{rybak_modelling_2006, mccrea_organization_2008, rybak_organization_2015}: each leg has a rhythm generator (RG) circuit, which produces the basic locomotor rhythm and controls its period, and a pattern formation (PF) circuit, which shapes the rhythm into phase-appropriate activation of the flexor and extensor motoneuron populations. A left--right rhythm generator circuit coordinates the two legs through commissural inhibition and excitation, following the models of left--right coordination \citep{molkov_mechanisms_2015, danner_central_2016}. The per-leg PF output drives the motoneuron pools --- here, the BPA activation commands for the hip, knee, and ankle muscle groups of each leg.

\subsection{Sensory Feedback Channels}

Three classes of proprioceptive feedback will close the loop between the MuJoCo plant and the CPG, mirroring the biological afferents organized in the Sensory Afferent Database \citep{bolen_sensory_2023}:

\begin{itemize}
    \item \textbf{Ia afferents} (muscle spindle primary endings) signal muscle length and velocity, providing stretch reflex pathways and phase-dependent modulation of the PF network.
    \item \textbf{Ib afferents} (Golgi tendon organs) signal muscle force, providing load-dependent regulation of stance-phase extensor activity.
    \item \textbf{II afferents} (muscle spindle secondary endings) signal length, contributing to position sensing and interjoint coordination.
    \item \textbf{Contact receptors} on the heel and toe provide discrete stance and swing boundary signals, driving phase transitions and weight-support switching.
\end{itemize}

\subsection{Verification Tests}

Before attempting locomotion, the network will be verified with simple simulation tests in which tonic stimulus is applied to the muscle motoneurons, to the pattern-formation layer, and to the rhythm-generator layer, and the resulting outputs are compared with the corresponding experimental literature. The closest biological analogue is the tonic-stimulation air-stepping protocol of Selionov et al. \citep{selionov_tonic_2009}, in which tonic central and sensory stimuli facilitate involuntary stepping in humans; analogous deletion and stimulation phenomena are reproduced by the two-level CPG model itself \citep{rybak_modelling_2006} and by neuromechanical implementations \citep{markin_afferent_2010, shevtsova_modeling_2016}. Comparing the simulated tonic-stimulation responses with these studies will validate the sign, gain, and phase structure of the implemented circuits before any walking-speed tuning is attempted.

\section{Extending the Sensorium: Vestibular, Ocular, and Cerebellar Contributions}\label{sec:fw_senses}

Once stepping is stable, three supraspinal contributions will be added. First, vestibular feedback will be simulated with IMU-like sensing of body orientation and angular velocity, extending the lab's balance-platform work, in which an inverted pendulum with pseudo-vestibular IMU feedback maintained balance on an oscillating platform \citep{bolen_balance_2024, bolen_selfbalancing_2024} and analytic neural networks reproduced human balance dynamics \citep{hilts_dynamic_2019}. Second, ocular feedback will be modeled to investigate how visual stabilization cues can help reset IMU drift --- a well-known practical problem for inertial sensing on robots --- by treating visual measurements of the horizon and surrounding as a low-frequency correction to the vestibular estimate. Third, a cerebellar layer will be explored as an adaptive gateway between higher-level locomotion commands and the CPG: a structure that can gate between different gaits or movement patterns (e.g., walking, running, turning) and tune the phase relationships among them, in line with the cerebellum's established role in adaptation and coordination. The proprioceptive--vestibular interaction literature \citep{akay_relative_2021} will guide how these signals are weighted.

\section{Physical Realization: A Treadmill Robot with Real-Time Neural Control}\label{sec:fw_hardware}

The final stage returns to hardware. The robot will walk on a treadmill with real-time feedback and control implemented on a embedded computing stack: Teensy microcontrollers for the low-latency per-joint loops (valve control, pressure sensing, contact sensing), and an Nvidia Jetson-class companion computer for the SNS network evaluation and activation computation. A serial connection to a host computer will be used only for data logging and high-level commands (e.g., run, walk, stimulate a named neural pathway), keeping the control loop fully embedded. This division of labor mirrors the SNS-Toolbox backend philosophy, which targets embedded computing platforms for real-time neural simulation \citep{nourse_sns_2023}.

Because the static force and torque models of Chapters~\ref{ch:methods}--\ref{ch:results} are computationally inexpensive lookup-table candidates, they will serve as the feedforward plant model inside the real-time controller, while the dynamic pressure--force behavior under pulse-modulated control will draw on the state-space model of \citet{elzein_pressure_2025}. Embedded execution of the complete CPG-plus-reflex network on this stack is the intended demonstration of the dissertation project: a robot that walks with human-like biomechanics under a biologically grounded neural controller, and whose nervous system --- unlike a living animal's --- can be instrumented, stimulated, and lesioned at will.

\section{Toward Dynamic Whole-Body Maneuvers}\label{sec:fw_kickflip}

Beyond steady walking, the long-term goal of this line of research is dynamic, whole-body maneuvering: transitions between gaits, reactive balancing, and ultimately ballistic maneuvers that require coordinating the full actuator set through brief, precisely timed activation bursts. The author's aspiration for this program is a robot that can perform a kickflip --- a maneuver that, like its skateboard namesake, demands explosive coordinated extension, precise timing, aerial body orientation, and controlled landing. It is exactly the kind of behavior that exercises every layer of the proposed architecture, from contact receptors and spinal pattern generation to vestibular orientation sensing and cerebellar coordination, and it would demonstrate that biomimetic actuation and neuromechanical control scale from steady locomotion to dynamic whole-body movement.
